\documentclass[11pt]{article}
\usepackage[margin=1in]{geometry}
\usepackage{amsmath, amssymb, booktabs}
\usepackage[hidelinks]{hyperref}

\title{Daily Residual Statistical Arbitrage in Recent Technology Equities\\
Replication}
\author{Thiago Amin}
\date{May 2026}

\begin{document}

\maketitle

\section{Objective and Scope}

Pairs trading begins with the idea that economically similar securities should
not diverge persistently after accounting for their common exposure. The
generalized version replaces a single peer stock with one or more factor
portfolios. A stock is bought when it appears abnormally cheap relative to its
factor-implied value and sold when it appears abnormally expensive, provided
the residual behavior is sufficiently consistent with mean reversion.

This replication asks whether that daily residual mean-reversion mechanism is
detectable in modern U.S. technology equities. The initial study will:
\begin{itemize}
    \item operate only at the daily trading frequency;
    \item use adjusted daily prices and traded ETF factors;
    \item estimate signals from trailing windows without future information;
    \item evaluate realized PNL after transaction costs; and
    \item report results for a recent fixed sample, proposed as January 2019
    through December 2025.
\end{itemize}

The first implementation is not intended to establish that daily trading is
the optimal horizon. It is a single-horizon baseline that can later become one
component of a multi-horizon allocation study.

\section{Data and Study Universe}

\subsection{Technology Stock Universe}

Let $\mathcal{U}_t$ denote the investable stock universe on trading day $t$.
The target universe is liquid U.S.-listed technology-sector equities. To avoid
survivorship bias in a formal replication, membership should be constructed
from point-in-time holdings of a technology benchmark ETF or another
point-in-time sector classification. A simplified exploratory implementation
may begin with a fixed set of liquid technology stocks, but results from that
variant must be labeled as conditional on the selected universe rather than as
a survivorship-free sector study.

For each stock $i\in\mathcal{U}_t$, let $S_{i,t}$ be its split- and
dividend-adjusted closing price. Daily log returns are
\[
    R_{i,t} = \log\left(\frac{S_{i,t}}{S_{i,t-1}}\right).
\]
Adjusted data are required because a mechanical dividend or split should not
be interpreted as an idiosyncratic pricing dislocation.

\subsection{Factor Benchmarks}

The primary factor specification follows the paper's interpretable ETF
approach. Because the trading universe is restricted to technology equities,
the primary factor is the traded technology-sector benchmark:
\[
    F_t = R_{\mathrm{XLK},t}.
\]
The stock is therefore assessed relative to the sector portfolio against which
it is economically comparable.

A pre-specified robustness analysis may use multiple traded benchmark returns,
for example
\[
    F_t =
    \begin{bmatrix}
        R_{\mathrm{SPY},t} &
        R_{\mathrm{XLK},t} &
        R_{\mathrm{QQQ},t}
    \end{bmatrix}^{\top}.
\]
This specification tests whether separating broad-market, technology-sector,
and technology-heavy growth exposure materially changes the residual signal.
Because these ETFs may be highly correlated, the single-factor XLK model
remains the primary clean replication benchmark.

\section{Residual Relative-Value Model}

\subsection{Factor Decomposition}

For one factor, the daily stock return is decomposed as
\[
    R_{i,t}
    =
    \alpha_i
    +
    \beta_i F_t
    +
    \widetilde{R}_{i,t}.
\]
The systematic component is $\beta_i F_t$ and
$\widetilde{R}_{i,t}$ is the idiosyncratic residual return. Under a
multi-factor specification,
\[
    R_{i,t}
    =
    \alpha_i
    +
    \sum_{j=1}^{m}\beta_{ij}F_{j,t}
    +
    \widetilde{R}_{i,t}.
\]
The research focus is not solely whether the factors explain return variance;
it is whether the remaining residual contains a stable and tradable
mean-reversion signal.

If a dollar exposure $Q_{i,t}$ is taken in stock $i$, the associated
factor-hedged position takes ETF exposures
\[
    Q_{F_j,t} = -\beta_{ij}Q_{i,t},
    \qquad j=1,\ldots,m.
\]
For that stock-level trade, estimated factor exposure is neutral:
\[
    \beta_{ij}Q_{i,t} + Q_{F_j,t} = 0.
\]
Aggregating these hedged stock trades produces a daily long--short portfolio
designed to earn residual convergence rather than broad technology-market
direction.

\subsection{Residual Level}

Let $M$ be the number of observations in the trailing estimation window. In
line with the source study, the baseline is
\[
    M = 60 \text{ business days}.
\]
At a decision date $t$, all parameters are estimated using observations ending
no later than $t-1$. Let $\widehat{\widetilde{R}}_{i,u}$ be the residual return
from the fitted factor regression for $u=t-M,\ldots,t-1$. Define the cumulative
residual level over the estimation window by
\[
    X_{i,u}
    =
    \sum_{v=t-M}^{u}\widehat{\widetilde{R}}_{i,v}.
\]
The residual level represents the cumulative relative-value displacement of
the stock from its factor hedge. A large positive value indicates accumulated
outperformance relative to the benchmark; a large negative value indicates
accumulated underperformance.

\section{Mean-Reversion Estimation}

\subsection{Ornstein--Uhlenbeck Specification}

For each stock, the residual level is modeled as an
Ornstein--Uhlenbeck process:
\[
    dX_i(u)
    =
    \kappa_i\left(m_i-X_i(u)\right)du
    +
    \sigma_i dW_i(u),
    \qquad \kappa_i>0.
\]
Here $m_i$ is the long-run residual equilibrium,
$\kappa_i$ is the speed of mean reversion, and $\sigma_i$ is residual
volatility. The conditional expected residual change is
\[
    \mathbb{E}_t[dX_i(u)]
    =
    \kappa_i\left(m_i-X_i(u)\right)du.
\]
Thus a residual above equilibrium predicts negative relative return, while a
residual below equilibrium predicts positive relative return.

\subsection{Discrete Daily Estimation}

With $\Delta=1/252$ year, the exact daily discretization can be estimated as
an AR(1) regression:
\[
    X_{i,u+\Delta}
    =
    a_i + b_i X_{i,u} + \eta_{i,u+\Delta}.
\]
The OU interpretation requires
\[
    0 < b_i < 1.
\]
For an accepted fit, the estimates map to mean-reversion quantities through
\[
    \widehat{\kappa}_i
    =
    -\frac{\log(\widehat{b}_i)}{\Delta},
    \qquad
    \widehat{m}_i
    =
    \frac{\widehat{a}_i}{1-\widehat{b}_i}.
\]
If $\widehat{\sigma}_{\eta,i}^{\,2}$ is the variance of the AR(1)
innovation, the estimated equilibrium standard deviation of the residual
level is
\[
    \widehat{\sigma}_{\mathrm{eq},i}
    =
    \frac{\widehat{\sigma}_{\eta,i}}
         {\sqrt{1-\widehat{b}_i^{\,2}}}.
\]

The estimated characteristic mean-reversion time is
\[
    \widehat{\tau}_i = \frac{1}{\widehat{\kappa}_i}.
\]
The source strategy restricts trading to residuals that revert sufficiently
quickly relative to the 60-day estimation window. The baseline eligibility
rule for this replication will therefore be
\[
    \widehat{\tau}_i < \frac{30}{252},
    \qquad \text{equivalently} \qquad
    \widehat{\kappa}_i > \frac{252}{30}.
\]
Stocks failing this rule, failing the stationary AR(1) restriction, or lacking
sufficient data are not eligible for a new position on that date.

\section{Trading Signal}

\subsection{S-Score}

The standardized residual displacement at decision date $t$ is
\[
    s_{i,t}
    =
    \frac{X_{i,t-1}-\widehat{m}_i}
         {\widehat{\sigma}_{\mathrm{eq},i}}.
\]
This score measures how far the stock's factor-relative residual lies from its
estimated equilibrium in equilibrium-standard-deviation units. The use of
$X_{i,t-1}$ reflects the implementation convention that a position for day
$t$ may use only information observable before that day's return is realized.

\subsection{Stateful Entry and Exit Rules}

The baseline daily signal retains the threshold structure reported in the
source study:
\[
\begin{array}{lll}
    \text{buy stock / sell hedge to open} & \text{if} & s_{i,t} < -1.25,\\
    \text{sell stock / buy hedge to open} & \text{if} & s_{i,t} > +1.25,\\
    \text{close short-stock position} & \text{if} & s_{i,t} < +0.75,\\
    \text{close long-stock position} & \text{if} & s_{i,t} > -0.50.
\end{array}
\]
The strategy is stateful: an entry signal opens a position and the
corresponding exit rule determines when it is unwound. It is intentionally a
discrete baseline rather than a continuously scaled function of the s-score,
so that the replication first tests the source trading mechanism.

\section{Portfolio Construction and PNL}

Let $z_{i,t}\in\{-1,0,+1\}$ denote the active stock-side signal for stock $i$
on day $t$, where $+1$ is long and $-1$ is short. Let $\Lambda_t$ be the
dollar amount assigned to each active stock-side position per unit of current
portfolio equity $E_t$. Stock and hedge dollars are
\[
    Q_{i,t}=E_t\Lambda_t z_{i,t},
    \qquad
    Q_{F_j,t}
    =
    -E_t\Lambda_t
    \sum_{i\in\mathcal{U}_t}z_{i,t}\widehat{\beta}_{ij,t}.
\]
The implementation will choose $\Lambda_t$ so that portfolio gross exposure
does not exceed a pre-specified budget. Both stock-side and ETF hedge
exposures must be included in this gross-exposure accounting.

Let $w_t$ denote the complete vector of stock and ETF portfolio weights
applied over day $t$, and let $R_t$ denote their realized daily return vector.
With transaction cost rate $c$, daily net PNL in return units is
\[
    \pi_t
    =
    w_t^{\top}R_t
    -
    c\left\|w_t-w_{t-1}\right\|_1.
\]
The timing convention is conservative: $w_t$ is fixed using information
available through day $t-1$ and earns only returns observed on day $t$. To
connect to the paper, the benchmark transaction-cost setting is
\[
    c = 0.0005,
\]
or five basis points per dollar of changed position. Results should also be
reported under lower and higher cost assumptions because modern liquidity and
implementation quality differ from the historical sample.

\section{Estimation and Backtest Protocol}

For each out-of-sample trading date $t$:
\begin{enumerate}
    \item Construct the eligible technology-stock universe using only
    point-in-time information available at $t$.
    \item Obtain adjusted daily returns for stocks and factor ETFs through
    date $t-1$.
    \item Using the trailing 60 business days, fit the factor regression for
    each stock and form its residual-level series.
    \item Fit the AR(1)/OU model and retain only stocks satisfying the
    stationarity and mean-reversion-time eligibility rules.
    \item Compute each eligible stock's s-score and update its open or closed
    position state.
    \item Construct the corresponding ETF hedges and rescale positions to the
    gross-exposure limit.
    \item Apply those positions to the return realized on date $t$, deduct
    turnover costs, and record net PNL.
\end{enumerate}

All model fitting and signal construction are rolling and out of sample. The
choice of universe, factors, 60-day window, thresholds, exposure budget, and
transaction-cost assumptions must be fixed before inspecting performance in
the test period. If alternative thresholds or factor sets are explored, they
will be reported as separate specifications rather than silently chosen from
the best backtest result.

\section{Results To Report}

The eventual implementation report should include:
\begin{itemize}
    \item universe construction and daily counts of eligible/traded stocks;
    \item factor specification and distributions of estimated factor loadings;
    \item distributions of $\widehat{\kappa}_i$,
    $\widehat{\tau}_i$, and s-scores;
    \item number of entries, exits, average holding period, and turnover;
    \item cumulative net PNL, annualized volatility, annualized Sharpe ratio,
    maximum drawdown, and hit rate;
    \item average long gross, short gross, hedge gross, net exposure, and
    factor exposure after hedging;
    \item transaction-cost sensitivity; and
    \item comparison of the primary XLK factor model against pre-specified
    factor robustness variants.
\end{itemize}

\section{Secondary Replication Extension: PCA Factors}

The original study also extracts systematic returns through principal
components. This is a useful second-stage replication, but it is not required
for the initial ETF implementation. For a trailing window, standardized stock
returns can be collected in a matrix $Y$ and the empirical correlation matrix
formed as
\[
    \rho = \frac{1}{M-1}YY^{\top}.
\]
If $v^{(1)},\ldots,v^{(k)}$ are eigenvectors corresponding to the retained
largest eigenvalues, their associated eigenportfolio returns become estimated
factors. The residual/OU/s-score/backtest procedure then remains unchanged;
only the source of the systematic factor returns changes. Comparing ETF and
PCA specifications within the same recent technology universe will test
whether the residual strategy is sensitive to how systematic technology risk
is extracted.

\section{Implementation Sequence}

\begin{enumerate}
    \item Implement adjusted daily price and ETF return ingestion.
    \item Implement point-in-time technology universe construction, with any
    simplified fixed-universe fallback clearly labeled.
    \item Implement rolling ETF regression residualization.
    \item Implement OU parameter estimation, eligibility screening, and
    s-score computation.
    \item Implement stateful factor-hedged positions, exposure limits, and
    transaction-cost-adjusted PNL.
    \item Produce the tables and plots specified above.
    \item Only after the daily replication is established, expose this strategy
    as a one-day sleeve for the endogenous-horizon study.
\end{enumerate}

\section*{Source Study}

Marco Avellaneda and Jeong-Hyun Lee,
\emph{Statistical Arbitrage in the U.S. Equities Market}, July 11, 2008,
SSRN abstract 1153505. 
\end{document}
